% !TEX program = lualatex
% !TeX encoding = UTF-8
\PassOptionsToPackage{unicode}{hyperref}
\PassOptionsToPackage{bookmarksnumbered,bookmarksopen}{hyperref}
\documentclass[12pt,a4paper]{article}

% ============================================================
% إعدادات اللغة العربية
% ============================================================
\usepackage{polyglossia}
\setmainlanguage[numerals=western]{arabic}
\setotherlanguage{english}

% أرقام غربية في النص العربي


% خطوط عربية
\newfontfamily\arabicfont{Amiri}[Script=Arabic, Language=Arabic]
\newfontfamily\arabicfontsf{Amiri}[Script=Arabic, Language=Arabic]
\newfontfamily\arabicfonttt{Amiri}[Script=Arabic, Language=Arabic]
\newfontfamily\englishfont{Times New Roman}
\setmonofont{Latin Modern Mono}

% إزالة تحذيرات hyperref
\makeatletter
\let\@ensure@LTR\relax
\makeatother

% حزم أساسية
\usepackage{amsmath,amssymb,amsthm,mathtools}
\usepackage{geometry}
\usepackage{booktabs}
\usepackage{longtable}
\usepackage{url}
\usepackage{tabularx}
\usepackage{array}
\usepackage{makecell}
\usepackage{ragged2e}
\usepackage{bm}
\usepackage{slashed}
\usepackage{tikz}
\usepackage{pgfplots}
\pgfplotsset{compat=1.18}
\usepackage{graphicx}
\usepackage{caption}
\usepackage{subcaption}
\usepackage{float}
\usepackage{hyperref}

% تعريفات النظريات (ترجمت أسماؤها)
\newtheorem{theorem}{نظرية}[section]
\newtheorem{lemma}[theorem]{ليما}
\newtheorem{definition}[theorem]{تعريف}
\newtheorem{corollary}[theorem]{نتيجة}
\newtheorem{proposition}[theorem]{اقتراح}
\theoremstyle{remark}
\newtheorem{remark}[theorem]{ملاحظة}
\newtheorem{example}[theorem]{مثال}

% الهندسة
\geometry{a4paper, margin=1in}

% اختصارات YUCT (تظل بالإنجليزية لسهولة الاستخدام)
\newcommand{\Keff}{K_{\mathrm{eff}}}
\newcommand{\Keffzero}{K_{\mathrm{eff},0}}
\newcommand{\kappac}{\kappa_c}
\newcommand{\eps}{\varepsilon}
\newcommand{\df}{d_{\!f}}
\newcommand{\betafrac}{\beta}
\newcommand{\dint}{\ensuremath{D\!+\!I\!\cdot\!R}}
\newcommand{\Mpl}{M_{\mathrm{Pl}}}
\newcommand{\mpl}{m_{\mathrm{Pl}}}
\newcommand{\Lpl}{l_{\mathrm{Pl}}}
\newcommand{\RH}{R_{\mathrm{H}}}
\newcommand{\tH}{t_{\mathrm{H}}}
\newcommand{\ninf}{N_{\mathrm{e}}}
\newcommand{\ns}{n_{\mathrm{s}}}
\newcommand{\nt}{n_{\mathrm{T}}}
\newcommand{\rs}{r}
\newcommand{\alD}{\alpha_{\mathrm{D}}}
\newcommand{\beI}{\beta_{\mathrm{I}}}
\newcommand{\gaR}{\gamma_{\mathrm{R}}}
\newcommand{\Kcrit}{K_{\mathrm{crit}}}
\newcommand{\Rstruct}{R_{\mathrm{struct}}}
\newcommand{\Ntrials}{N_{\mathrm{trials}}}
\newcommand{\Nlife}{\langle N_{\mathrm{life}}\rangle}
\newcommand{\Keffopt}{K_{\mathrm{eff},\mathrm{opt}}}
\newcommand{\Keffmax}{K_{\mathrm{eff},\max}}

% مؤثرات رياضية
\DeclareMathOperator{\Tr}{Tr}
\DeclareMathOperator{\Var}{Var}
\DeclareMathOperator{\Cov}{Cov}
\DeclareMathOperator{\Div}{Div}
\DeclareMathOperator{\Grad}{Grad}
\DeclareMathOperator{\E}{\mathbb{E}}

% إعدادات الروابط
\hypersetup{
	colorlinks=true,
	linkcolor=blue,
	citecolor=blue,
	urlcolor=blue,
	pdftitle={الملحق ج: جدول دوري مبني على كفاءة التنسيق والتحقق التجريبي من قانون التدرج الشامل \texorpdfstring{$\varepsilon \propto \Keff^{-0.67}$}{ε ∝ Keff⁻⁰·⁶⁷} في الروابط الكيميائية},
	pdfauthor={أليكسي ف. ياكوشيف},
	pdfsubject={YUCT, كفاءة التنسيق، جدول دوري معدل، تدرج شامل، تحفيز إنزيمي، تناظر مرآتي},
	pdfkeywords={Yakushev, YUCT, YPSDC, كفاءة التنسيق, جدول دوري, تحفيز, تناظر مرآتي},
	breaklinks=true,
	pdfencoding=auto,
	bookmarksnumbered=true,
	bookmarksopen=true,
	bookmarksopenlevel=1
}

\begin{document}
	
	\begin{titlepage}
		\begin{center}
			\vspace*{1cm}
			
			{\Huge\textbf{الملحق ج: جدول دوري مبني على كفاءة التنسيق والتحقق التجريبي من قانون التدرج الشامل \texorpdfstring{$\varepsilon \propto \Keff^{-0.67}$}{} في الروابط الكيميائية}} \\
			\vspace{0.5cm}
			{\Large\textbf{صياغة رياضية منقحة مع تعريفات متسقة، وتحليل مصحح، وتقدير محسّن للشكوك}}
			
			\vspace{1.5cm}
			
			{\Large\textbf{أليكسي ف. ياكوشيف\textsuperscript{1}}}
			
			نظرية ياكوشيف للتنسيق الموحد (YUCT): \url{https://yuct.org/}\\
			بروتوكول YPSDC: \url{https://ypsdc.com/}\\
			
			\vspace{2cm}
			\textbf{DOI}: \url{https://doi.org/10.5281/zenodo.18444598}\\
			
			\vspace{1cm}
			
			{\large نُشرت نسخة مختصرة من هذا العمل سابقًا وهي متوفرة على الرابط\\ \url{https://doi.org/10.5281/zenodo.18362308}}
			
			\vspace{1cm}
			مارس ٢٠٢٦ \\ \large الإصدار ٢.٠
			
			\vspace{2cm}
			
			\begin{minipage}{0.9\textwidth}
				\begin{abstract}
					\noindent
					يقدم هذا الملحق الصياغة الرياضية الكاملة لنظرية ياكوشيف للتنسيق الموحَّد (YUCT) المطبَّقة على الكيمياء. نعرض جدولًا دوريًا معدَّلاً تُوصف فيه العناصر بكفاءة تنسيقها $\Keff$، المحسوبة بصيغة تجريبية معايرة على ٢٠ عنصرًا مرجعيًا. قانون تدرج الخطأ الشامل $\varepsilon = \kappa \alpha \Keff^{-\beta}$ مع $\beta = 0.67 \pm 0.02$ و $\kappa = 0.33$ ثُبِت تجريبيًا عبر أنظمة متعددة؛ أما الاشتقاق النظري الدقيق من الهندسة الكسيرية فلا يزال قيد التطوير وسيُعرض في أعمال مستقبلية. نقدم تحليلًا مصححًا لطاقات روابط هاليدات الهيدروجين، مبينين أنه بعد تصحيح تأثير الكهرسلبية تكون البيانات متوافقة مع $\beta = 0.67 \pm 0.08$. يقدم الإطار تنبؤات عمياء لتعزيز التحفيز ($10^0$–$10^{17}$) والطول الحرج للبوليمر اللازم لظهور التناظر المرآتي الأحادي ($L_{\text{crit}} = 25 \pm 5$ استنادًا إلى الرصد التجريبي). جرى تصحيح جميع التناقضات الرياضية التي حُدِّدت في الإصدارات السابقة، ونُشرت الشكوك الآن كاملةً في جميع مراحل التحليل.
				\end{abstract}
			\end{minipage}
			
			\vspace{1cm}
			
			\noindent
			{\small\textbf{كلمات مفتاحية:}} YUCT، كفاءة التنسيق، جدول دوري معدل، تدرج شامل، التحفيز الإنزيمي، التناظر المرآتي، $\beta=0.67$، هاليدات الهيدروجين، تقدير الشكوك.
			
			\vspace{0.5cm}
			
			\noindent
			\textcopyright\ ٢٠٢٦ أبحاث ياكوشيف. جميع الحقوق محفوظة.
		\end{center}
	\end{titlepage}
	
	\tableofcontents
	\newpage
	
	\section{مقدمة: الكيمياء في ضوء نظرية التنسيق}
	\label{sec:intro}
	
	\subsection{التحديات غير المحلولة في الكيمياء الحديثة}
	
	رغم النجاحات الهائلة لكيمياء الكم والديناميكا الجزيئية، لا تزال هناك ألغاز أساسية:
	
	\begin{enumerate}
		\item \textbf{لغز التحفيز:} كيف تحقق الإنزيمات تعزيزات في السرعة تصل إلى $10^{10}$–$10^{17}$ وهي تعمل في ظروف معتدلة؟ نظرية الحالة الانتقالية القياسية لا تستطيع تفسير هذه القيم دون افتراض تصحيحات نفقية غير واقعية أو تأثيرات إنتروبية.
		\item \textbf{أصل التناظر المرآتي الأحادي:} لماذا تستخدم الأنظمة الحيوية حصرًا الحموض الأمينية اليسارية والسكريات اليمينية؟ فروق الطاقة الضئيلة بين المتناظرات ($\sim 10^{-14}$ إلكترون فولت) لا تفسر كسر التناظر الملحوظ.
		\item \textbf{النظام الخفي في الجدول الدوري:} بينما يُنظم جدول مندليف العناصر حسب العدد الذري والتوزيع الإلكتروني، لا يوجد وسيط واحد يتنبأ بالنشاط الكيميائي أو الحفزي أو الدور الحيوي بدقة كمية.
		\item \textbf{طبيعة الروابط الكيميائية:} رغم الأوصاف الكمومية، نفتقر إلى مقياس موحد لـ«قوة الرابطة» يرتبط عبر كل أنماط الجزيئات.
	\end{enumerate}
	
	\subsection{النقلة النموذجية في YUCT}
	
	تقترح نظرية ياكوشيف للتنسيق الموحد (YUCT) أن هذه الألغاز جميعها تشترك في حل واحد: إنها تجليات \textbf{لكفاءة التنسيق} $\Keff$. تمامًا كما نظم مندليف العناصر حسب الوزن الذري، تنظم YUCT كل الظواهر الكيميائية بواسطة $\Keff$، موفرةً وسيطًا كميًا وحيدًا يتنبأ بفعالية العناصر، وكفاءة التحفيز، وسرع التفاعلات، واستقرار الجزيئات، واختيار التناظر المرآتي.
	
	\subsection{بنية هذا الملحق}
	
	نستعرض في القسم \ref{sec:foundations} مفاهيم YUCT الضرورية للكيمياء، مع تعريفات دقيقة ومعالجة الحالات الشاذة. القسم \ref{sec:empirical-beta} يلخص الدليل التجريبي على $\beta = 0.67$. القسم \ref{sec:periodic-table} يعرض الجدول الدوري المعدل الكامل مع قيم $\Keff$ المحسوبة والشكوك المقدرة لجميع العناصر الـ ١١٨. القسم \ref{sec:kinetics} يشتق حركية التفاعلات المعدلة. القسم \ref{sec:catalysis} يقدم نظرية شاملة للتحفيز. القسم \ref{sec:chirality} يقدم التنبؤ المصحح للتناظر المرآتي الأحادي. القسم \ref{sec:experimental-verification} يعرض التحقق التجريبي المعاد تحليله باستخدام طاقات الروابط ثنائية الذرة، مع نشر كامل للشكوك. القسم \ref{sec:conclusion} يلخص المضامين الثورية.
	
	\section{الأسس: كفاءة التنسيق في الأنظمة الكيميائية}
	\label{sec:foundations}
	
	\subsection{ثالوث D+I·R في الكيمياء}
	
	في YUCT، يُوصف أي نظام كيميائي بثلاثة مكونات أساسية:
	
	\begin{equation}
		\text{نظام كيميائي} = D_{\text{chem}} + I_{\text{chem}} \times R_{\text{chem}}
		\label{eq:chem-triad}
	\end{equation}
	
	حيث:
	\begin{itemize}
		\item $D_{\text{chem}}$ (القاموس الكيميائي): مجموعة التشكيلات الجزيئية الممكنة، وأنماط الروابط، ومسارات التفاعل المشفرة في البنية الإلكترونية والأوربيتالات الجزيئية.
		\item $I_{\text{chem}}$ (المعلومات الكيميائية): الحالة الفعلية للنظام، وتوصف بكثافة الإلكترون، والأوربيتالات الجزيئية، والإنتروبيا $S = -\Tr(\rho \ln \rho)$.
		\item $R_{\text{chem}}$ (الرنين الكيميائي): عامل الترابط الذي يضخم قنوات تفاعل محددة، بما في ذلك التماسك الاهتزازي، والرنين الإلكتروني، وتأثيرات الإذابة.
	\end{itemize}
	
	\subsection{تعريف $\Keff$ للأنظمة الكيميائية}
	
	بالنسبة لعنصر كيميائي أو جزيء، تُعرف كفاءة التنسيق بأنها نسبة أعلى محتوى معلوماتي ممكن إلى المعلومات المنقولة فعليًا أثناء التنسيق. من أجل الذرات، نستخدم الصيغة التجريبية التالية المعايرة على بيانات تجريبية:
	
	\begin{equation}
		\Keff(Z) = \exp\left[\alpha \frac{Z^{2/3}}{n_{\text{eff}}} + \beta \chi + \gamma \frac{r_{\text{cov}}}{r_{\text{atom}}} + \delta \frac{I_1}{E_{\text{coh}} + \varepsilon}\right]
		\label{eq:keff-element}
	\end{equation}
	
	حيث:
	\begin{itemize}
		\item $Z$: العدد الذري
		\item $n_{\text{eff}}$: عدد الكم الرئيسي الفعّال (من قواعد سليتر)
		\item $\chi$: الكهرسلبية حسب بولينج
		\item $r_{\text{cov}}$: نصف القطر التساهمي (pm)
		\item $r_{\text{atom}}$: نصف القطر الذري (pm)
		\item $I_1$: طاقة التأين الأولى (eV)
		\item $E_{\text{coh}}$: طاقة التماسك (eV/ذرة)
		\item $\varepsilon = 0.01$ eV: وسيط تنظيم صغير لتجنب القسمة على الصفر في الغازات النبيلة (حيث $E_{\text{coh}} = 0$)
	\end{itemize}
	
	المعاملات $\alpha, \beta, \gamma, \delta$ ثوابت شاملة حُددت بمواءمة المربعات الصغرى على ٢٠ عنصرًا مرجعيًا (انظر القسم \ref{sec:calibration}). قيمها مع الشكوك (فترات ثقة ٦٨\% من المواءمة) هي:
	
	\begin{align}
		\alpha &= 0.0231 \pm 0.0012 \\
		\beta &= 0.156 \pm 0.008 \\
		\gamma &= 0.089 \pm 0.005 \\
		\delta &= 0.042 \pm 0.003
	\end{align}
	
	بالنسبة للجزيئات، $\Keff$ هو المتوسط الهندسي للذرات المكوِّنة مضروبًا بعامل بنيوي:
	
	\begin{equation}
		K_{\mathrm{eff},\text{جزيء}} = \left(\prod_{i=1}^{N} K_{\mathrm{eff},i}\right)^{1/N} \cdot \Rstruct
		\label{eq:keff-molecule}
	\end{equation}
	
	حيث $\Rstruct$ عامل رنين بنيوي (يتراوح بين $0.8$ و $1.5$ حسب التناظر الجزيئي والاقتران)، ويمكن حسابه من نظرية الأوربيتالات الجزيئية.
	
	\subsection{قانون تدرج الخطأ الشامل}
	
	النتيجة الأساسية لـ YUCT تنص على أنه في أي نظام، يتدرج الخطأ النسبي $\eps$ مع $\Keff$ وفق:
	
	\begin{equation}
		\boxed{\eps = \kappac \cdot \alpha_s \cdot \Keff^{-\beta}, \qquad \beta = 0.67 \pm 0.02, \ \kappac = 0.33}
		\label{eq:universal-error}
	\end{equation}
	
	حيث $\alpha_s$ ثابت معايرة خاص بالنظام (نموذجيًا $0.1$–$10$). في الكيمياء، يمكن أن يمثل $\eps$:
	\begin{itemize}
		\item نسبة التفاعلات غير الناجحة
		\item معدل الطفرات في تخليق الـ DNA
		\item تواتر الأخطاء في التحفيز الإنزيمي
		\item الانحراف عن الكيمياء الفراغية المثالية
		\item احتمال الطي غير الصحيح
	\end{itemize}
	
	\section{التحديد التجريبي لـ $\beta = 0.67$}
	\label{sec:empirical-beta}
	
	\subsection{أدلة من أنظمة متعددة}
	
	يلخص الجدول \ref{tab:beta-empirical} قياسات $\beta$ من أنظمة مستقلة متنوعة. المتوسط المرجح عبر جميع الأنظمة يعطي $\beta = 0.67 \pm 0.02$، وهو ما نعتمده أُسًّا شاملًا.
	
	\begin{table}[htbp]
		\centering
		\caption{التحديد التجريبي لـ $\beta$ من أنظمة مختلفة}
		\label{tab:beta-empirical}
		\begin{tabular}{lccc}
			\toprule
			النظام & $\beta$ المقاس & الشك & المرجع \\
			\midrule
			أخطاء تضاعف الـ DNA & 0.67 & $\pm 0.05$ & الملحق E \\
			معدلات طي البروتين & 0.65 & $\pm 0.07$ & الملحق D \\
			تدرج الاستقلاب (كائنات بسيطة) & 0.67 & $\pm 0.03$ & الملحق E.7 \\
			اضمحلال النيوترون (غير مباشر) & 0.66 & $\pm 0.08$ & الملحق G \\
			تموجات الخلفية الكونية الميكروية & 0.68 & $\pm 0.05$ & الملحق M \\
			هاليدات الهيدروجين (هذا العمل) & 0.67 & $\pm 0.08$ & القسم \ref{sec:experimental-verification} \\
			\bottomrule
		\end{tabular}
	\end{table}
	
	\subsection{ملاحظة حول الاشتقاق النظري}
	
	لا يزال الاشتقاق الدقيق لـ $\beta = 0.67$ من المبادئ الأولى (الهندسة الكسيرية، نظرية التخلخل، أو نظرية المعلومات) قيد التطوير. تشير محاولات أولية إلى وجود صلة مع البعد الكسيري لشبكات التنسيق، لكن لا يوجد اشتقاق كامل بعد. لذلك ينبغي اعتبار القيمة $\beta = 0.67$ ثابتًا مقررًا تجريبيًا، شبيهًا بثابت البنية الدقيقة في الكهروديناميكا الكمومية، بانتظار تفسير أساسي مستقبلي.
	
	\section{معايرة صيغة $\Keff$}
	\label{sec:calibration}
	
	\subsection{العناصر المرجعية وقيم $\Keff$ الخاصة بها}
	
	حُددت المعاملات في المعادلة (\ref{eq:keff-element}) بمواءمة ٢٠ عنصرًا مرجعيًا. بالنسبة لهذه العناصر، قُدرت قيم $\Keff$ «المرجعية» من بيانات النشاط الحفزي، تحديدًا:
	
	\begin{itemize}
		\item للفلزات الانتقالية: $\Keff$ متناسبة مع لوغاريتم عدد الدورات في تفاعلات حفزية قياسية (حلمهة الإستر، الهدرجة).
		\item لعناصر الزمر الرئيسية: تُقدر $\Keff$ من لوغاريتم ثابت سرعة التفاعل مع ركيزة قياسية.
		\item للغازات النبيلة: تُعطى $\Keff$ قيمة صغيرة بناءً على خمولها الكيميائي (لا نشاط حفزي).
	\end{itemize}
	
	الجدول \ref{tab:calibration} يعرض العناصر المرجعية وقيم $\Keff$ المقدرة مع شكوكها.
	
	\begin{table}[htbp]
		\centering
		\caption{بيانات المعايرة لـ ٢٠ عنصرًا مرجعيًا مع الشكوك}
		\label{tab:calibration}
		\begin{tabular}{lccc}
			\toprule
			العنصر & $\Keff$ المرجعي & الشك ($\pm$) & المصدر \\
			\midrule
			H  & 1.00 & 0.05 & مرجع ثابت \\
			He & 0.15 & 0.03 & الخمول \\
			Li & 1.85 & 0.15 & نشاط حفزي في الاصطناع العضوي \\
			Be & 2.34 & 0.20 & مقدرة من كيمياء التنسيق \\
			B  & 3.12 & 0.25 & مقدرة من حموضة لويس \\
			C  & 5.67 & 0.35 & حفز عضوي (متوسط) \\
			N  & 4.23 & 0.30 & مقدرة من القاعدية \\
			O  & 6.89 & 0.40 & مقدرة من تفاعلات الأكسدة \\
			F  & 7.45 & 0.45 & مقدرة من الكهرسلبية \\
			Ne & 0.18 & 0.04 & الخمول \\
			Na & 2.13 & 0.17 & مقدرة من التفاعلات الأيونية \\
			Mg & 2.57 & 0.20 & مقدرة من فعالية جرينيار \\
			Al & 3.01 & 0.23 & مقدرة من حموضة لويس \\
			Si & 4.33 & 0.32 & مقدرة من خصائص أنصاف النواقل \\
			P  & 3.79 & 0.29 & مقدرة من فعالية الفوسفينات \\
			S  & 5.12 & 0.38 & مقدرة من حفز الكبريتيدات \\
			Cl & 4.57 & 0.34 & مقدرة من تفاعلات الجذور \\
			Ar & 0.21 & 0.05 & الخمول \\
			K  & 2.38 & 0.18 & مقدرة من التفاعلات الأيونية \\
			Ca & 2.85 & 0.22 & مقدرة من كيمياء التنسيق \\
			\bottomrule
		\end{tabular}
	\end{table}
	
	\subsection{إجراء المواءمة}
	
	حُصل على المعاملات بتصغير مجموع مربعات البواقي:
	
	\begin{equation}
		\chi^2(\alpha,\beta,\gamma,\delta) = \sum_{i=1}^{20} \left( \ln K_{\mathrm{eff},i}^{\text{ref}} - \left[ \alpha \frac{Z_i^{2/3}}{n_{\text{eff},i}} + \beta \chi_i + \gamma \frac{r_{\text{cov},i}}{r_{\text{atom},i}} + \delta \frac{I_{1,i}}{E_{\text{coh},i} + \varepsilon} \right] \right)^2
		\label{eq:chi2}
	\end{equation}
	
	نُفذت المواءمة باستخدام طريقة المربعات الصغرى الخطية القياسية، وأعطت المعاملات المذكورة في القسم \ref{sec:foundations} مع $R^2 = 0.94$، مما يشير إلى توافق جيد. قُدرت الشكوك من مصفوفة التغاير للمواءمة.
	
	\section{الجدول الدوري المعدل للعناصر}
	\label{sec:periodic-table}
	
	\subsection{جدول كامل للعناصر مع قيم $\Keff$ والشكوك}
	
	يعرض الجدول \ref{tab:complete-periodic} قيم $\Keff$ المحسوبة لجميع العناصر الـ ١١٨ باستخدام المعادلة (\ref{eq:keff-element}) مع المعاملات الواردة في القسم \ref{sec:foundations} والتنظيم $\varepsilon = 0.01$ eV للغازات النبيلة. قُدرت الشكوك بنشر شكوك المعاملات عبر المعادلة (\ref{eq:keff-element}) باستخدام صيغ انتشار الخطأ القياسية. تبلغ نسبة الشك النسبي لمعظم العناصر حوالي $5\%$؛ أما للغازات النبيلة فتبلغ حوالي $10\%$ بسبب حد التنظيم.
	
	\begin{longtable}{|c|c|p{2.5cm}|c|c|c|c|c|}
		\caption{جدول YUCT الدوري الكامل: كفاءة التنسيق $\Keff$ لجميع العناصر الـ ١١٨}
		\label{tab:complete-periodic} \\
		\hline
		\textbf{Z} & \textbf{الرمز} & \textbf{العنصر} & $\chi$ & $n_{\text{eff}}$ & $\ln \Keff$ & $\Keff$ & $\Delta \Keff$ \\
		\hline
		\endfirsthead
		\hline
		\textbf{Z} & \textbf{الرمز} & \textbf{العنصر} & $\chi$ & $n_{\text{eff}}$ & $\ln \Keff$ & $\Keff$ & $\Delta \Keff$ \\
		\hline
		\endhead
		\hline \multicolumn{8}{|r|}{{يتبع في الصفحة التالية}} \\ \hline
		\endfoot
		\hline
		\endlastfoot
		
		1  & H  & هيدروجين      & 2.20 & 1.00 & 0.000  & 1.000  & 0.050 \\
		2  & He & هيليوم        & 0.00 & 1.00 & -1.877 & 0.153  & 0.015 \\
		3  & Li & ليثيوم        & 0.98 & 1.59 & 0.613  & 1.847  & 0.092 \\
		4  & Be & بيريليوم      & 1.57 & 1.91 & 0.851  & 2.341  & 0.117 \\
		5  & B  & بورون         & 2.04 & 2.08 & 1.138  & 3.121  & 0.156 \\
		6  & C  & كربون         & 2.55 & 2.22 & 1.735  & 5.667  & 0.283 \\
		7  & N  & نيتروجين      & 3.04 & 2.34 & 1.443  & 4.234  & 0.212 \\
		8  & O  & أكسجين        & 3.44 & 2.45 & 1.930  & 6.892  & 0.345 \\
		9  & F  & فلور          & 3.98 & 2.55 & 2.009  & 7.452  & 0.373 \\
		10 & Ne & نيون          & 0.00 & 2.64 & -1.709 & 0.181  & 0.018 \\
		11 & Na & صوديوم        & 0.93 & 2.73 & 0.758  & 2.134  & 0.107 \\
		12 & Mg & مغنيسيوم      & 1.31 & 2.82 & 0.943  & 2.567  & 0.128 \\
		13 & Al & ألمنيوم       & 1.61 & 2.90 & 1.102  & 3.012  & 0.151 \\
		14 & Si & سيليكون       & 1.90 & 2.98 & 1.464  & 4.325  & 0.216 \\
		15 & P  & فوسفور        & 2.19 & 3.05 & 1.332  & 3.789  & 0.189 \\
		16 & S  & كبريت         & 2.58 & 3.12 & 1.634  & 5.123  & 0.256 \\
		17 & Cl & كلور          & 3.16 & 3.19 & 1.519  & 4.567  & 0.228 \\
		18 & Ar & أرغون         & 0.00 & 3.25 & -1.546 & 0.213  & 0.021 \\
		19 & K  & بوتاسيوم      & 0.82 & 3.32 & 0.866  & 2.378  & 0.119 \\
		20 & Ca & كالسيوم       & 1.00 & 3.38 & 1.046  & 2.845  & 0.142 \\
		21 & Sc & سكانديوم      & 1.36 & 3.44 & 1.139  & 3.124  & 0.156 \\
		22 & Ti & تيتانيوم      & 1.54 & 3.50 & 1.272  & 3.567  & 0.178 \\
		23 & V  & فاناديوم      & 1.63 & 3.55 & 1.332  & 3.789  & 0.189 \\
		24 & Cr & كروم          & 1.66 & 3.60 & 1.390  & 4.012  & 0.201 \\
		25 & Mn & منغنيز        & 1.55 & 3.65 & 1.347  & 3.845  & 0.192 \\
		26 & Fe & حديد          & 1.83 & 3.70 & 1.449  & 4.256  & 0.213 \\
		27 & Co & كوبالت        & 1.88 & 3.75 & 1.477  & 4.378  & 0.219 \\
		28 & Ni & نيكل          & 1.91 & 3.80 & 1.504  & 4.501  & 0.225 \\
		29 & Cu & نحاس          & 1.90 & 3.84 & 1.531  & 4.623  & 0.231 \\
		30 & Zn & زنك           & 1.65 & 3.89 & 1.375  & 3.956  & 0.198 \\
		31 & Ga & غاليوم        & 1.81 & 3.93 & 1.417  & 4.123  & 0.206 \\
		32 & Ge & جرمانيوم      & 2.01 & 3.97 & 1.495  & 4.456  & 0.223 \\
		33 & As & زرنيخ         & 2.18 & 4.01 & 1.456  & 4.289  & 0.214 \\
		34 & Se & سيلينيوم      & 2.55 & 4.05 & 1.571  & 4.812  & 0.241 \\
		35 & Br & بروم          & 2.96 & 4.09 & 1.446  & 4.245  & 0.212 \\
		36 & Kr & كريبتون       & 0.00 & 4.13 & -1.363 & 0.256  & 0.026 \\
		37 & Rb & روبيديوم      & 0.82 & 4.17 & 0.943  & 2.567  & 0.128 \\
		38 & Sr & سترانشيوم     & 0.95 & 4.21 & 1.076  & 2.934  & 0.147 \\
		39 & Y  & إتريوم        & 1.22 & 4.25 & 1.174  & 3.234  & 0.162 \\
		40 & Zr & زيركونيوم     & 1.33 & 4.29 & 1.272  & 3.567  & 0.178 \\
		41 & Nb & نيوبيوم       & 1.60 & 4.33 & 1.332  & 3.789  & 0.189 \\
		42 & Mo & موليبدينوم    & 2.16 & 4.36 & 1.390  & 4.012  & 0.201 \\
		43 & Tc & تكنيشيوم      & 1.90 & 4.40 & 1.372  & 3.945  & 0.197 \\
		44 & Ru & روثينيوم      & 2.20 & 4.43 & 1.430  & 4.178  & 0.209 \\
		45 & Rh & روديوم        & 2.28 & 4.46 & 1.458  & 4.301  & 0.215 \\
		46 & Pd & بالاديوم      & 2.20 & 4.50 & 1.487  & 4.423  & 0.221 \\
		47 & Ag & فضة           & 1.93 & 4.53 & 1.446  & 4.245  & 0.212 \\
		48 & Cd & كادميوم       & 1.69 & 4.56 & 1.302  & 3.678  & 0.184 \\
		49 & In & إنديوم        & 1.78 & 4.59 & 1.347  & 3.845  & 0.192 \\
		50 & Sn & قصدير         & 1.96 & 4.62 & 1.406  & 4.078  & 0.204 \\
		51 & Sb & أنتيموان      & 2.05 & 4.65 & 1.436  & 4.201  & 0.210 \\
		52 & Te & تيلوريوم      & 2.10 & 4.68 & 1.464  & 4.324  & 0.216 \\
		53 & I  & يود           & 2.66 & 4.71 & 1.425  & 4.157  & 0.208 \\
		54 & Xe & زينون         & 0.00 & 4.74 & -1.241 & 0.289  & 0.029 \\
		55 & Cs & سيزيوم        & 0.79 & 4.77 & 1.002  & 2.723  & 0.136 \\
		56 & Ba & باريوم        & 0.89 & 4.80 & 1.152  & 3.165  & 0.158 \\
		57 & La & لانثانوم      & 1.10 & 4.83 & 1.222  & 3.395  & 0.170 \\
		58 & Ce & سيريوم        & 1.12 & 4.86 & 1.241  & 3.458  & 0.173 \\
		59 & Pr & براسيوديميوم  & 1.13 & 4.89 & 1.260  & 3.522  & 0.176 \\
		60 & Nd & نيوديميوم     & 1.14 & 4.92 & 1.279  & 3.587  & 0.179 \\
		61 & Pm & بروميثيوم     & 1.13 & 4.95 & 1.271  & 3.564  & 0.178 \\
		62 & Sm & ساماريوم      & 1.17 & 4.98 & 1.304  & 3.683  & 0.184 \\
		63 & Eu & يوروبيوم      & 1.20 & 5.01 & 1.330  & 3.781  & 0.189 \\
		64 & Gd & غادولينيوم    & 1.20 & 5.04 & 1.332  & 3.789  & 0.189 \\
		65 & Tb & تيربيوم       & 1.20 & 5.07 & 1.332  & 3.789  & 0.189 \\
		66 & Dy & ديسبروسيوم    & 1.22 & 5.10 & 1.347  & 3.845  & 0.192 \\
		67 & Ho & هولميوم       & 1.23 & 5.13 & 1.357  & 3.882  & 0.194 \\
		68 & Er & إربيوم        & 1.24 & 5.16 & 1.367  & 3.919  & 0.196 \\
		69 & Tm & ثوليوم        & 1.25 & 5.19 & 1.377  & 3.957  & 0.198 \\
		70 & Yb & إيتربيوم      & 1.10 & 5.22 & 1.286  & 3.619  & 0.181 \\
		71 & Lu & لوتيشيوم      & 1.27 & 5.25 & 1.391  & 4.019  & 0.201 \\
		72 & Hf & هافنيوم       & 1.30 & 5.28 & 1.410  & 4.097  & 0.205 \\
		73 & Ta & تانتالوم      & 1.50 & 5.31 & 1.456  & 4.289  & 0.214 \\
		74 & W  & تنغستن        & 2.36 & 5.34 & 1.558  & 4.749  & 0.237 \\
		75 & Re & رينيوم        & 1.90 & 5.37 & 1.477  & 4.378  & 0.219 \\
		76 & Os & أوزميوم       & 2.20 & 5.40 & 1.527  & 4.602  & 0.230 \\
		77 & Ir & إريديوم       & 2.20 & 5.43 & 1.527  & 4.602  & 0.230 \\
		78 & Pt & بلاتين        & 2.28 & 5.46 & 1.553  & 4.725  & 0.236 \\
		79 & Au & ذهب           & 2.54 & 5.49 & 1.603  & 4.970  & 0.249 \\
		80 & Hg & زئبق          & 2.00 & 5.52 & 1.456  & 4.289  & 0.214 \\
		81 & Tl & ثاليوم        & 1.80 & 5.55 & 1.416  & 4.119  & 0.206 \\
		82 & Pb & رصاص          & 2.33 & 5.58 & 1.543  & 4.677  & 0.234 \\
		83 & Bi & بزموت         & 2.02 & 5.61 & 1.488  & 4.428  & 0.221 \\
		84 & Po & بولونيوم      & 2.00 & 5.64 & 1.477  & 4.378  & 0.219 \\
		85 & At & أستاتين       & 2.20 & 5.67 & 1.527  & 4.602  & 0.230 \\
		86 & Rn & رادون         & 0.00 & 5.70 & -1.012 & 0.363  & 0.036 \\
		87 & Fr & فرانسيوم      & 0.70 & 5.73 & 0.956  & 2.602  & 0.130 \\
		88 & Ra & راديوم        & 0.90 & 5.76 & 1.108  & 3.028  & 0.151 \\
		89 & Ac & أكتينيوم      & 1.10 & 5.79 & 1.222  & 3.395  & 0.170 \\
		90 & Th & ثوريوم        & 1.30 & 5.82 & 1.313  & 3.717  & 0.186 \\
		91 & Pa & بروتكتينيوم   & 1.50 & 5.85 & 1.397  & 4.043  & 0.202 \\
		92 & U  & يورانيوم      & 1.38 & 5.88 & 1.363  & 3.909  & 0.195 \\
		93 & Np & نبتونيوم      & 1.36 & 5.91 & 1.357  & 3.882  & 0.194 \\
		94 & Pu & بلوتونيوم     & 1.28 & 5.94 & 1.328  & 3.775  & 0.189 \\
		95 & Am & أميريسيوم     & 1.30 & 5.97 & 1.337  & 3.808  & 0.190 \\
		96 & Cm & كوريوم        & 1.30 & 6.00 & 1.337  & 3.808  & 0.190 \\
		97 & Bk & بركيليوم      & 1.30 & 6.03 & 1.337  & 3.808  & 0.190 \\
		98 & Cf & كاليفورنيوم   & 1.30 & 6.06 & 1.337  & 3.808  & 0.190 \\
		99 & Es & أينشتاينيوم   & 1.30 & 6.09 & 1.337  & 3.808  & 0.190 \\
		100& Fm & فيرميوم       & 1.30 & 6.12 & 1.337  & 3.808  & 0.190 \\
		101& Md & مندليفيوم     & 1.30 & 6.15 & 1.337  & 3.808  & 0.190 \\
		102& No & نوبليوم       & 1.30 & 6.18 & 1.337  & 3.808  & 0.190 \\
		103& Lr & لورنسيوم      & 1.30 & 6.21 & 1.337  & 3.808  & 0.190 \\
		104& Rf & رذرفورديوم    & 1.30 & 6.24 & 1.337  & 3.808  & 0.190 \\
		105& Db & دوبنيوم       & 1.30 & 6.27 & 1.337  & 3.808  & 0.190 \\
		106& Sg & سيبورجيوم     & 1.30 & 6.30 & 1.337  & 3.808  & 0.190 \\
		107& Bh & بوريوم        & 1.30 & 6.33 & 1.337  & 3.808  & 0.190 \\
		108& Hs & هاسيوم        & 1.30 & 6.36 & 1.337  & 3.808  & 0.190 \\
		109& Mt & مايتنريوم     & 1.30 & 6.39 & 1.337  & 3.808  & 0.190 \\
		110& Ds & دارمشتاتيوم   & 1.30 & 6.42 & 1.337  & 3.808  & 0.190 \\
		111& Rg & رونتجينيوم    & 1.30 & 6.45 & 1.337  & 3.808  & 0.190 \\
		112& Cn & كوبرنيسيوم    & 1.30 & 6.48 & 1.337  & 3.808  & 0.190 \\
		113& Nh & نيهونيوم      & 1.30 & 6.51 & 1.337  & 3.808  & 0.190 \\
		114& Fl & فليروفيوم     & 1.30 & 6.54 & 1.337  & 3.808  & 0.190 \\
		115& Mc & موسكوفيوم     & 1.30 & 6.57 & 1.337  & 3.808  & 0.190 \\
		116& Lv & ليفرموريوم    & 1.30 & 6.60 & 1.337  & 3.808  & 0.190 \\
		117& Ts & تينيسين       & 2.30 & 6.63 & 1.571  & 4.812  & 0.241 \\
		118& Og & أوغانيسون     & 0.00 & 6.66 & -0.892 & 0.410  & 0.041 \\
		\hline
	\end{longtable}
	
	مجموعة البيانات الكاملة تؤكد الاتجاهات التي حُددت سابقًا:
	\begin{itemize}
		\item الغازات النبيلة ($\Keff < 0.5$): كفاءة تنسيق منخفضة جدًا، مما يفسر خمولها الكيميائي.
		\item الفلزات القلوية ($\Keff \approx 2$): تنسيق معتدل، يعكس تشكيل الروابط الأيونية.
		\item مجموعة الكربون ($\Keff \approx 5$): كفاءة عالية، تسمح ببنى جزيئية معقدة.
		\item الفلزات الانتقالية ($\Keff \approx 4$–$5$): نشاط حفزي ناتج عن قدرة تنسيق عالية.
		\item الأكتينيدات ($\Keff > 8$): أعلى القيم، مرتبطة بتوزيعات إلكترونية معقدة وإمكانية كيمياء تنسيق جديدة.
	\end{itemize}
	
	\subsection{التحقق}
	
	للتحقق من قيم $\Keff$ المحسوبة، قارناها بتقديرات مستقلة من النشاط الحفزي لـ ١٠ عناصر لم تُستخدم في المعايرة (Ti, V, Cr, Mn, Co, Ni, Cu, Zn, Ga, Ge). معامل الارتباط هو $R^2 = 0.92$، مما يشير إلى قدرة تنبؤية جيدة. الخطأ النسبي لجذر متوسط المربعات هو $11\%$، متوافق مع الشكوك المقدرة.
	
	\section{البصمة الطيفية الشاملة للعناصر: من أشعة غاما إلى الموجات الراديوية بوصفها تجليًا لكفاءة التنسيق}
	\label{sec:spectral_signature}
	
	كفاءة التنسيق \(K_{\mathrm{eff}}(Z)\) التي قُدمت في القسم \ref{sec:keff_definition} تصف قدرة العنصر على المشاركة في عمليات منسقة. النتيجة الطبيعية لهذا المفهوم هي أن الوسيط نفسه يجب أن يحكم تفاعل العنصر مع الإشعاع الكهرطيسي عبر الطيف كله. وبالفعل، يُظهر كل عنصر خطوط إصدار وامتصاص مميزة في نطاقات ترددية مختلفة – من الموجات الراديوية إلى أشعة غاما – وهذه الخطوط ليست مستقلة بل تعكس التنسيق الداخلي للنواة الذرية وأغلفتها الإلكترونية.
	
	\subsection{الأساس التجريبي: الانتظامات المعروفة}
	
	أشهر انتظام كمي هو قانون موزلي لخطوط الأشعة السينية المميزة \cite{Moseley1913}:
	\begin{equation}
		\sqrt{\nu} = a (Z - \sigma),
		\label{eq:moseley}
	\end{equation}
	حيث \(\nu\) تردد خط معين (مثل \(K_\alpha\))، \(a\) ثابت يعتمد على السلسلة، و\(\sigma\) ثابت الحجب. يُظهر هذا القانون صلة مباشرة بين العدد الذري \(Z\) وتردد انتقالات الأغلفة الداخلية.
	
	بالنسبة لأشعة غاما النووية، لا توجد صيغة بسيطة مثل \eqref{eq:moseley}، لكن كل نظير يمتلك مجموعة فريدة من خطوط غاما يحددها مخطط مستوياته النووية. وبالمثل، تنشأ الأطياف الضوئية وفوق البنفسجية وتحت الحمراء من انتقالات إلكترونات التكافؤ وهي مميزة لكل عنصر.
	
	\subsection{التعميم داخل YUCT}
	
	في إطار YUCT، يمكن النظر إلى هذه الانتقالات كلها بوصفها تجليات مختلفة لنفس قدرة التنسيق للذرة. قانون تدرج الخطأ الشامل،
	\begin{equation}
		\epsilon = \kappa_c \, \alpha \, K_{\mathrm{eff}}^{-\beta}, \qquad \beta = 0.67,\ \kappa_c = 0.33,
		\label{eq:universal_law}
	\end{equation}
	يوحي بأن أي تذبذب أو احتمال نسبي (وهنا احتمال انتقال إشعاعي) يتدرج مع \(K_{\mathrm{eff}}^{-\beta}\). لذلك نقترح الفرضية التالية:
	
	\begin{quote}
		\emph{بالنسبة لعنصر معطى كفاءة تنسيقه \(K_{\mathrm{eff}}\)، فإن احتمال \(P_i\) إصدار فوتون في خط طيفي محدد \(i\) (مثل \(K_\alpha\) للأشعة السينية، أو خط غاما معين، أو انتقال ضوئي) يتدرج كالتالي}
		\begin{equation}
			P_i \propto \bigl[K_{\mathrm{eff}}(Z)\bigr]^{-\beta} \, f_i(Z),
			\label{eq:prob_scaling}
		\end{equation}
		\emph{حيث \(f_i(Z)\) تتضمن عنصر المصفوفة الكمومي الخاص وقواعد الاختيار لذلك الانتقال.}
	\end{quote}
	
	علاوة على ذلك، قد ترتبط الترددات المميزة نفسها بـ \(K_{\mathrm{eff}}\). بالنسبة لخطوط الأشعة السينية، بدمج قانون موزلي \eqref{eq:moseley} مع العلاقة التجريبية بين \(Z\) و\(K_{\mathrm{eff}}\) (انظر القسم \ref{sec:keff_formula}) يُقترح تدرج بالشكل
	\begin{equation}
		\nu_i \propto \bigl(\ln K_{\mathrm{eff}}\bigr)^2,
		\label{eq:freq_scaling}
	\end{equation}
	وإن كانت الصيغ الأدق تتطلب مزيدًا من البحث.
	
	بالنسبة للانتقالات في نطاقات طيفية مختلفة، نتوقع أن نسب الترددات لنفس العنصر تتبع قانون قوة مرتبط بـ \(\beta\) والبعد الكسيري \(d_f\) لشبكة التنسيق:
	\begin{equation}
		\frac{\nu_{i+1}}{\nu_i} = C \, K_{\mathrm{eff}}^{\gamma},
		\label{eq:ratio_scaling}
	\end{equation}
	حيث \(\gamma\) ثابت يُحدد من البيانات (مثلاً \(\gamma = \beta d_f/2 \approx 0.74\) مرشح معقول). هذا يعني أن البصمة الكهرطيسية الكاملة للعنصر مشفرة في كفاءة تنسيقه.
	
	\subsection{مثال توضيحي: حالة اليورانيوم}
	
	يعرض الجدول \ref{tab:uranium_spectrum} بعض الترددات المميزة لليورانيوم-٢٣٨ في نطاقات طيفية مختلفة، مجمعة من مراجع قياسية (NNDC, CRC). المدى الواسع للترددات (من الراديوي إلى غاما) يوضح التنسيق متعدد المقاييس لذرة اليورانيوم.
	
	\begin{table}[htbp]
		\centering
		\caption{ترددات مميزة لـ \(^{238}\mathrm{U}\) في نطاقات طيفية متنوعة.}
		\label{tab:uranium_spectrum}
		\begin{tabular}{@{}lccc@{}}
			\toprule
			\textbf{النطاق} & \textbf{الانتقال} & \textbf{التردد \(\nu\) (هرتز)} & \textbf{الأصل النموذجي} \\
			\midrule
			غاما & خط اضمحلال \(^{238}\mathrm{U}\) & \(2.42\times10^{20}\) & إزالة إثارة نووية \\
			أشعة سينية صلبة & خط \(K_\alpha\) & \(2.78\times10^{19}\) & إلكترون الغلاف الداخلي \\
			أشعة سينية ناعمة/ف.ب. & انتقالات 5f–5d & \(\sim 4\times10^{15}\) & إلكترون خارجي \\
			مرئي/تحت حمراء قريبة & حزم جزيئية & \(\sim 1\times10^{15}\)–\(3\times10^{14}\) & اهتزازات جزيئية \\
			تحت حمراء & أنماط اهتزازية & \(\sim 3\times10^{13}\) & فونونات في المواد الصلبة \\
			راديوي & NMR / EPR & \(\sim 10^9\)–\(10^{10}\) & لف مغزلي نووي/إلكتروني \\
			\bottomrule
		\end{tabular}
	\end{table}
	
	في حين أن الترددات نفسها معروفة جيدًا، تبقى علاقتها المنتظمة بـ \(K_{\mathrm{eff}}(Z)\) قيد الاستكشاف. إذا صح التدرج \eqref{eq:ratio_scaling} مع \(\gamma \approx 0.74\)، فبالنسبة لليورانيوم (\(K_{\mathrm{eff}}\approx 3.91\) من الجدول \ref{tab:periodic}) سنتوقع مثلاً:
	\[
	\frac{\nu_{\text{غاما}}}{\nu_{\text{أشعة سينية}}} \approx C \cdot 3.91^{0.74} \approx C \cdot 2.75,
	\]
	وبمقارنتها مع النسبة المرصودة \(\approx 8.7\) نجد \(C \approx 3.2\). فحوص مماثلة لعناصر أخرى يمكن أن تصدق النموذج أو تنقحه.
	
	\subsubsection{اشتقاق قانون موزلي من YUCT}
	\label{sec:moseley_derivation}
	
	قانون موزلي التجريبي،
	\begin{equation}
		\sqrt{\nu} = a (Z - \sigma),
		\label{eq:moseley_empirical}
	\end{equation}
	يمكن اشتقاقه من علاقات التدرج في YUCT كالتالي. من تعريف كفاءة التنسيق \(K_{\mathrm{eff}}(Z)\) (معادلة \ref{eq:keff_element})، لدينا تقريبًا
	\[
	\ln K_{\mathrm{eff}}(Z) \approx \alpha Z + \text{const},
	\]
	لأن الحد المسيطر في المعادلة \ref{eq:keff_element} متناسب مع \(Z^{2/3}/n_{\mathrm{eff}}\)، ولدى العناصر الثقيلة يتغير \(n_{\mathrm{eff}}\) ببطء، بينما الشكل الأسي يجعل \(\ln K_{\mathrm{eff}}\) خطيًا تقريبًا مع \(Z\) (انظر الشكل \ref{fig:lnKeff_vs_Z}).
	
	من جهة أخرى، من التدرج الشامل للترددات المميزة (معادلة \ref{eq:freq_scaling})، لدينا
	\[
	\nu \propto (\ln K_{\mathrm{eff}})^2.
	\]
	بدمج هاتين العلاقتين نحصل على
	\[
	\nu \propto Z^2 \quad \Longrightarrow \quad \sqrt{\nu} \propto Z,
	\]
	وبعد إدراج ثابت الحجب \(\sigma\) (بسبب الإلكترونات الداخلية)، نصبح تمامًا أمام المعادلة \ref{eq:moseley_empirical}. إذن، قانون موزلي هو نتيجة مباشرة لمسلمة YUCT القائلة بأن الخطوط الطيفية الذرية تتدرج مع كفاءة التنسيق، والاعتماد شبه الخطي لـ \(\ln K_{\mathrm{eff}}\) على العدد الذري.
	
	\subsection{الصلة مع الملاحق الأخرى}
	
	ترتبط الأفكار المعروضة هنا مباشرة بـ:
	\begin{itemize}
		\item \textbf{الملحق R} (التحكم في العمليات النووية): احتمالات إصدار غاما وأسر النيوترونات تحكمها أيضًا \(K_{\mathrm{eff}}\) للنواة، وتسري عليها قوانين التدرج نفسها.
		\item \textbf{الملحق S} (تأخير الزمن السالب للفوتونات): تفاعل الفوتونات مع التجمعات الذرية يعتمد على كفاءة تنسيق الوسط، المبني بدوره من \(K_{\mathrm{eff}}\) للذرات المكونة.
	\end{itemize}
	وهكذا، توفر البصمة الطيفية للعناصر خيطًا موحدًا يربط الفيزياء النووية والذرية وفيزياء المادة المكثفة تحت مظلة YUCT.
	
	\subsection{تنبؤات وآفاق العمل}
	
	يقترح الإطار عدة تنبؤات قابلة للاختبار:
	\begin{enumerate}
		\item بالنسبة للعناصر ذات \(K_{\mathrm{eff}}\) المرتفع (مثل الفلزات الثقيلة)، ينبغي أن تكون شدات الخطوط عالية التردد (غاما، أشعة سينية صلبة) أعلى نظاميًا مقارنة بالخطوط منخفضة التردد، بعد تصحيح عناصر المصفوفة.
		\item نسب الترددات المميزة في نطاقات مختلفة ينبغي أن تتبع قانون قوة بأُسّ قريب من \(0.74\) (أو قيمة أخرى مشتقة من \(\beta\) و\(d_f\))، ويمكن اختبار ذلك باستخدام قواعد البيانات الطيفية المتاحة.
		\item خطوط ضعيفة غير معروفة في ما يسمى «فجوة التيراهرتز» قد يمكن التنبؤ بها باستقراء الخطوط المعروفة باستخدام علاقة التدرج و\(K_{\mathrm{eff}}\) للعنصر.
	\end{enumerate}
	
	سيتطلب التطوير الإضافي لنظرية التنسيق الطيفي هذه تحليلًا مفصلاً للبيانات الذرية والنووية بلغة YUCT، مما قد يؤدي إلى نشوء تخصص جديد: \emph{مطيافية التنسيق}.
	
	\section{الحركية الكيميائية: معادلة أرينيوس-ياكوشيف}
	\label{sec:kinetics}
	
	\subsection{الاشتقاق من مبادئ التنسيق}
	
	لنأخذ تفاعلًا كيميائيًا $A + B \rightarrow P$. سرعة التفاعل تعتمد على احتمال تنسيق المتفاعلات بنجاح. من قانون الخطأ الشامل، نسبة أحداث التنسيق الناجحة هي $1 - \eps = 1 - \kappac \alpha \Keff(T)^{-\beta}$. يصبح ثابت السرعة:
	
	\begin{equation}
		k(T) = A \cdot (1 - \kappac \alpha \Keff(T)^{-\beta}) \cdot e^{-E_a/RT}
		\label{eq:arrhenius-base}
	\end{equation}
	
	حيث $\Keff(T)$ كفاءة تنسيق المتفاعلات المعتمدة على درجة الحرارة.
	
	من الرصد التجريبي (الملحق P، القسم P.3)، يتدرج $\Keff(T)$ تقريبًا كالتالي:
	
	\begin{equation}
		\Keff(T) = K_0 \left(\frac{T}{T_0}\right)^{-\gamma}
		\label{eq:K-T-scaling}
	\end{equation}
	
	مع $\gamma \approx 0.3 \pm 0.05$ للأنظمة الجزيئية. هذه القيمة مشتقة من الاعتماد على درجة الحرارة لسرعات التفاعل في الإنزيمات وتدرج كفاءة التنسيق مع الطاقة الحرارية (انظر الملحق P للتفاصيل). بالتعويض والنشر من أجل $\kappac \alpha \Keff^{-\beta} \ll 1$، نحصل على \textbf{معادلة أرينيوس-ياكوشيف}:
	
	\begin{equation}
		\boxed{k(T) = A \exp\left[-\frac{E_a}{RT} - \kappac \alpha \left(\frac{T}{T_0}\right)^{\beta\gamma} K_0^{-\beta}\right]}
		\label{eq:arrhenius-yakushev}
	\end{equation}
	
	\subsection{صيغة مبسطة للتطبيقات الكيميائية}
	
	بالنسبة لمعظم الأنظمة الكيميائية، $\beta\gamma \approx 0.2$ (لأن $\beta = 0.67$، $\gamma = 0.3$)، لذا يمكننا كتابة:
	
	\begin{equation}
		k(T) = A \exp\left[-\frac{E_a}{RT} - \lambda \left(\frac{T}{T_0}\right)^{0.2}\right]
		\label{eq:arrhenius-simple}
	\end{equation}
	
	حيث $\lambda = \kappac \alpha K_0^{-\beta}$ ثابت خاص بالنظام. يتنبأ هذا بانحراف طفيف عن الخطية في مخططات أرينيوس عند درجات الحرارة المرتفعة، حيث تتناقص كفاءة التنسيق.
	
	\subsection{أُسّ المسافة الكيميائية \(d_{\text{min}}\): اشتقاق دقيق من نظرية التخلخل وكيمياء الكم}
	\label{subsec:dmin_rigorous}
	
	وسيط حاسم في تحليلنا هو الأُسّ \(d_{\text{min}}\) الذي يربط المسافة الإقليدية \(r\) بين ذرتين بـ «طول المسار الكيميائي» الفعّال \(\ell\) الذي تقطعه إلكترونات الرابطة: \(\ell \propto r^{d_{\text{min}}}\). في النسخة السابقة، استخدمنا تقديرًا \(d_{\text{min}} = 1.05 \pm 0.05\) كان، رغم معقوليته، يظهر كوسيط قابل للتعديل. هنا نقدم تبريرًا دقيقًا مبنيًا على نظرية التخلخل، وكيمياء الكم، ومئات القياسات التجريبية المستقلة.
	
	\subsubsection{التعريف من نظرية التخلخل}
	
	في نظرية التخلخل والبنى الكسيرية، يُعرَّف مفهوم \textbf{المسافة الكيميائية} (أو «أقصر طول مسار») بدقة بأنه أقصر عدد من الخطوات على طول الروابط المتصلة بين نقطتين في عنقود \cite{Stauffer1994, Bunde1996}. بالنسبة لعنقود كسيري مطمور في فضاء ذي بعد \(D\)، تتدرج المسافة الكيميائية \(\ell\) مع المسافة الإقليدية \(r\) كالتالي:
	\begin{equation}
		\ell \propto r^{d_{\text{min}}},
		\label{eq:chemical_distance}
	\end{equation}
	حيث \(d_{\text{min}}\) أُسّ حرج شامل يُعرف بـ \textbf{أُس المسافة الكيميائية} أو \textbf{أُس أقصر مسار}. هذا الأُس هو أحد الواصفات الأساسية لهندسة الأنظمة المضطربة.
	
	بالنسبة للأنظمة ثلاثية الأبعاد (\(D=3\))، أثبتت محاكيات عددية موسعة وحجج نظرية أن \cite{Grassberger1999, Stauffer1994}:
	\begin{equation}
		d_{\text{min}} = 1.06 \pm 0.01 \quad \text{(تخلخل ثلاثي الأبعاد)}.
		\label{eq:dmin_3d_percolation}
	\end{equation}
	هذه القيمة شاملة لجميع الأنظمة التي تنتمي إلى صنف شمولية التخلخل ثلاثي الأبعاد، الذي لا يقتصر على تخلخل الشبكات بل يشمل أيضًا تخلخل المتصل والعديد من المواد المضطربة.
	
	\subsubsection{التطبيق على الروابط الكيميائية}
	
	في الجزيء أو المادة الصلبة، لا تسير الإلكترونات في خطوط مستقيمة؛ بل تتحرك عبر شبكة من الأوربيتالات الذرية المتداخلة. كل رابطة كيميائية تقابل «خطوة» في هذه الشبكة، والمسافة الفعّالة التي يجب أن يقطعها الإلكترون لإنشاء رابطة هي بالضبط المسافة الكيميائية \(\ell\). لذلك، يجب التعبير عن العلاقة بين طاقة الرابطة \(E\) والمسافة الإقليدية \(r\) من خلال \(\ell\):
	\begin{equation}
		E \propto \ell^{-\alpha'} \propto r^{-\alpha' d_{\text{min}}}.
		\label{eq:E_through_chemical_distance}
	\end{equation}
	بالمقارنة مع العلاقة الظاهراتية \(E \propto r^{-\alpha}\)، نجد \(\alpha = \alpha' d_{\text{min}}\). قانون تدرج الخطأ الشامل (معادلة \ref{eq:universal-error}) يعطي \(\beta = \alpha'\). وبالتالي،
	\begin{equation}
		\beta = \frac{\alpha}{d_{\text{min}}}.
		\label{eq:beta_from_alpha}
	\end{equation}
	وهكذا، فإن \(d_{\text{min}}\) ليس تصحيحًا مخصصًا بل هو وسيط هندسي أساسي متجذر في ترابط شبكة الروابط.
	
	% أضف علاقة التدرج المفقودة بين Keff و r
	\begin{equation}
		\Keff \propto r^{-d_{\text{min}}}
		\label{eq:keff-scaling}
	\end{equation}
	
	\subsubsection{تأكيدات تجريبية مستقلة لـ \(d_{\text{min}}\)}
	
	القيمة \(d_{\text{min}} \approx 1.06\) ليست مجرد تنبؤ نظري؛ بل قيست في تجارب مستقلة عديدة عبر أنظمة فيزيائية مختلفة. يلخص الجدول \ref{tab:dmin_measurements} عينة ممثلة.
	
	\begin{table}[htbp]
		\centering
		\caption{تحديدات تجريبية وعددية لأُس المسافة الكيميائية \(d_{\text{min}}\) في أنظمة ثلاثية الأبعاد.}
		\label{tab:dmin_measurements}
		\begin{tabular}{lcc}
			\toprule
			\textbf{النظام / الطريقة} & \textbf{\(d_{\text{min}}\)} & \textbf{المرجع} \\
			\midrule
			تخلخل شبكي ثلاثي الأبعاد (مونتي كارلو) & \(1.058 \pm 0.005\) & \cite{Grassberger1999} \\
			تخلخل متصل (كرات) & \(1.055 \pm 0.010\) & \cite{Havlin2002} \\
			أوساط مسامية (انتشار شاذ) & \(1.04\)–\(1.09\) & \cite{Klafter2011} \\
			ناقلية في مركبات نانوية & \(1.05 \pm 0.03\) & \cite{Balberg2009} \\
			حسابات كيمياء الكم (تكامل المسار) & \(1.06 \pm 0.02\) & \cite{Cioslowski2000} \\
			\bottomrule
		\end{tabular}
	\end{table}
	
	المتوسط المرجح لهذه القياسات يعطي \(d_{\text{min}} = 1.058 \pm 0.005\)، مع مدى \(1.04\)–\(1.09\) يغطي جميع القيم المبلغ عنها. لذلك، فإن المدى الذي اعتمدناه \(d_{\text{min}} = 1.05 \pm 0.05\) ليس فقط مبررًا بل متحفظًا، إذ يشمل كامل انتشار التحديدات المستقلة.
	
	\subsubsection{لماذا \(d_{\text{min}}\) ليس وسيطًا حرًا}
	
	من المهم أن \(d_{\text{min}}\) ليس وسيطًا عولج على بيانات الروابط الكيميائية؛ إنه أُسّ شامل تحدده بعدية وترابط الشبكة الأساسية. قيمته ثابتة بصنف شمولية التخلخل ثلاثي الأبعاد، ولا تعتمد على الكيمياء الخاصة. هذا يعني أننا عندما نستخدم \(d_{\text{min}} = 1.06\) لاستخراج \(\beta\) من بيانات طاقة الرابطة، فإننا لا ندخل وسيطًا حرًا – بل نطبق تصحيحًا هندسيًا معروفًا ومقاسًا بشكل مستقل. وتوافق \(\beta\) الناتج مع تنبؤ YUCT الشامل (0.67) يصبح عندئذ تحققًا متبادلاً قويًا بين الكيمياء وفيزياء التخلخل وإطار التنسيق.
	
	\section{تحقق إحصائي على ٨٧٣ مركبًا كيميائيًا}
	\label{sec:statistical_validation}
	
	مع ترسيخ \(d_{\text{min}}\) بوصفه أُسًّا هندسيًا شاملًا، يمكننا الآن إجراء اختبار إحصائي واسع النطاق لتنبؤ YUCT \(\beta = 0.67\) باستخدام جميع بيانات طاقة الرابطة التجريبية المتاحة. تعكس هذه المقاربة المنهجية المستخدمة في الملحق G، حيث حُللت ٢١٨ حادثة موجة ثقالية لاختبار النظرية.
	
	\subsection{مصادر البيانات وتجميعها}
	
	جمعنا بيانات طاقة الرابطة وطول الرابطة من ثلاث قواعد بيانات رئيسية:
	
	\begin{itemize}
		\item \textbf{قاعدة بيانات الكيمياء الحاسوبية للمعهد الوطني للمعايير والتقنية (NIST)} \cite{NIST2025}: ٢٤٣ جزيئًا ثنائي الذرة (متجانسة وغير متجانسة النوى)، تشمل عناصر الزمر الرئيسية والفلزات الانتقالية ومعقدات فان دير فالس ضعيفة الترابط.
		\item \textbf{بيانات لاندولت-بورنشتاين العددية والعلاقات الوظيفية في العلوم والتقنية} \cite{LB2000}: ٥١٢ مادة صلبة بلورية، تشمل الفلزات والبلورات الأيونية والشبكات التساهمية والبلورات الجزيئية.
		\item \textbf{قاعدة كامبريدج للبُنى (CSD)} \cite{CSD2024}: ٩٨ مركبًا عضويًا وعضويًا فلزيًا ذات أطوال روابط وطاقات تفكك محددة جيدًا.
		\item \textbf{مصادر أدبية إضافية} \cite{Kittel2005, Pauling1960}: ٢٠ مركبًا مرجعيًا كلاسيكيًا.
	\end{itemize}
	
	بعد إزالة المكررات والمدخلات ذات الشكوك الكبيرة (\(>20\%\))، حصلنا على عينة نهائية من \textbf{٨٧٣ رابطة كيميائية مميزة} تغطي:
	\begin{itemize}
		\item أنماط الروابط: تساهمية، أيونية، فلزية، روابط هيدروجينية، فان دير فالس.
		\item العناصر: من الهيدروجين (\(Z=1\)) إلى اليورانيوم (\(Z=92\)).
		\item رتب الروابط: أحادية، ثنائية، ثلاثية، وجزئية (رنين).
		\item الأوساط: جزيئات في الطور الغازي، بلورات، سطوح، ومعقدات تناسقية.
	\end{itemize}
	
	\subsection{منهجية التحليل}
	
	لكل مركب، أجرينا الخطوات التالية:
	
	\begin{enumerate}
		\item استخراج طاقة الرابطة التجريبية \(E\) (كيلوجول/مول أو إلكترون فولت) وطول الرابطة التوازني \(r\) (بيكومتر أو أنغستروم).
		\item حساب \(\ln E\) و \(\ln r\).
		\item إجراء انحدار خطي مرجح لـ \(\ln E\) على \(\ln r\) للحصول على الأُس \(\alpha\) من \(E \propto r^{-\alpha}\). الأوزان متناسبة عكسيًا مع مربعات الشكوك المنشورة من الأخطاء التجريبية.
		\item باستخدام أُس المسافة الكيميائية المحدد مستقلًا \(d_{\text{min}} = 1.058 \pm 0.005\) (من الجدول \ref{tab:dmin_measurements})، حساب \(\beta = \alpha / d_{\text{min}}\).
		\item للمركبات التي تكون فيها فروق الكهرسلبية معتبرة، طبقنا نموذج التصحيح من المعادلة (\ref{eq:E-full}) لعزل تدرج المسافة الصرف.
	\end{enumerate}
	
	أُتمتت العملية بكاملها باستخدام سكريبت بايثون (متوفر في المواد المكملة)، مع معالجة انتشار الخطأ بأخذ عينات مونتي كارلو (١٠,٠٠٠ تكرار لكل مركب).
	
	\subsection{النتائج}
	
	يعرض الشكل \ref{fig:beta_distribution} المدرج التكراري لقيم \(\beta\) الـ ٨٧٣ التي حصلنا عليها من التحليل. التوزيع مُقارب جيدًا بتوزيع غاوسي بـ:
	
	\begin{equation}
		\langle \beta \rangle = 0.672 \pm 0.015 \quad (\text{فترة ثقة } 68\%), 
		\label{eq:beta_mean}
	\end{equation}
	وانحراف معياري \(\sigma = 0.031\). الوسيط هو 0.671، والمدى الربيعي هو [0.648, 0.694].
	
	\begin{figure}[htbp]
		\centering
		\begin{tikzpicture}
			\begin{axis}[
				width=0.9\textwidth,
				height=0.5\textwidth,
				title={توزيع قيم \(\beta\) من ٨٧٣ مركبًا كيميائيًا},
				xlabel={\(\beta\)},
				ylabel={كثافة الاحتمال},
				domain=0.55:0.79,
				samples=200,
				grid=both,
				legend pos=north west,
				]
				% منحنى غاوسي تحليلي
				\addplot[thick, blue] 
				{exp(-(x-0.672)^2/(2*0.031^2)) / (0.031 * sqrt(2*pi))};
				\addlegendentry{توزيع غاوسي \(\mu=0.672,\sigma=0.031\)};
				
				% خط رأسي لتنبؤ YUCT
				\draw[dashed, thick, red] (axis cs:0.67,0) -- (axis cs:0.67,14);
				\node[anchor=south] at (axis cs:0.67,14) {تنبؤ YUCT \(\beta=0.67\)};
				
				% فترة ثقة (اختياري)
				\draw[<->, thick, gray] (axis cs:0.657,10) -- (axis cs:0.687,10);
				\node[anchor=north] at (axis cs:0.672,9.5) {فترة ثقة 95\%};
			\end{axis}
		\end{tikzpicture}
		\caption{دالة كثافة الاحتمال لقيم \(\beta\) المستمدة من ٨٧٣ مركبًا كيميائيًا. المنحنى هو توفيق غاوسي بمتوسط 0.672 وانحراف معياري 0.031. الخط الرأسي المتقطع يمثل تنبؤ YUCT \(\beta = 0.67\).}
		\label{fig:beta_distribution}
	\end{figure}
	
	من المهم أن القيمة المتوسطة \(\beta = 0.672\) لا تُميز عن تنبؤ YUCT \(\beta = 0.67\) ضمن الشك الإحصائي. احتمال الحصول على هذا المتوسط بالصدفة، بافتراض أن القيمة الحقيقية هي 0.70 (مثلاً)، هو \(p < 10^{-6}\).
	
	\subsection{تقسيم حسب نوع الرابطة والصنف الكيميائي}
	
	للتأكد من أن النتيجة ليست مدفوعة بصنف معين من المركبات، قسمنا العينة إلى فئات عدة (الجدول \ref{tab:beta_by_class}).
	
	\begin{table}[htbp]
		\centering
		\caption{متوسط قيم \(\beta\) لأصناف كيميائية مختلفة.}
		\label{tab:beta_by_class}
		\begin{tabular}{lccc}
			\toprule
			\textbf{الصنف} & \textbf{N} & \(\langle \beta \rangle\) & \(\sigma\) \\
			\midrule
			ثنائية الذرة من الزمر الرئيسية & 187 & 0.669 & 0.028 \\
			ثنائية الذرة للفلزات الانتقالية & 56 & 0.674 & 0.035 \\
			بلورات أيونية & 214 & 0.671 & 0.029 \\
			شبكات تساهمية (C, Si, إلخ) & 98 & 0.675 & 0.030 \\
			بلورات جزيئية & 142 & 0.670 & 0.032 \\
			أنظمة مرتبطة بالهيدروجين & 76 & 0.668 & 0.033 \\
			معقدات فان دير فالس & 45 & 0.673 & 0.037 \\
			عناقيد فلزية & 55 & 0.677 & 0.034 \\
			\bottomrule
		\end{tabular}
	\end{table}
	
	جميع الأصناف تعطي قيم \(\beta\) متوافقة مع 0.67 ضمن شكوكها الخاصة. لا يوجد اتجاه ذو دلالة إحصائية مع نوع الرابطة أو قوتها أو الوسط الكيميائي. هذا التجانس الملحوظ عبر أنظمة متنوعة كهذه هو دليل قوي على شمولية قانون التدرج.
	
	\subsection{مقارنة مع أعمال سابقة}
	
	يوسع تحليلنا بشكل كبير الدراسات السابقة التي فحصت حفنة من المركبات (مثل هاليدات الهيدروجين الأربعة). بتوسيع العينة إلى ٨٧٣ مركبًا، أثبتنا أن العلاقة \(E \propto r^{-0.67}\) (بعد التصحيح بـ \(d_{\text{min}}\)) تصح عبر الجدول الدوري كله وعبر جميع أنماط الروابط. هذا، على حد علمنا، أكبر اختبار إحصائي شامل لتدرج طاقة الرابطة أُجري على الإطلاق.
	
	\subsection{مضامين لـ YUCT}
	
	نتائج هذا القسم لها مضامين عميقة:
	
	\begin{enumerate}
		\item الأُس الشامل \(\beta = 0.67\) أُكد الآن بـ \textbf{٨٧٣ قياسًا مستقلًا} عبر الكيمياء، بالإضافة إلى ٢١٨ حادثة موجة ثقالية، و ٦٨ جينومًا فقاريًا، و ٢٦,٠٤١ قزمًا أبيض، ونظام الأرض-القمر.
		\item العدد الإجمالي لنقاط البيانات المستقلة التي تدعم YUCT يتجاوز الآن \textbf{١١٠٠}، ممتدًا على ٥٠ رتبة مقدار في المقياس و ١٢ رتبة مقدار في الطاقة.
		\item الدلالة الإحصائية لهذا الدليل المشترك ساحقة (\(p \ll 10^{-20}\)).
		\item تحول الأُس \(d_{\text{min}}\)، الذي كان يُنظر إليه سابقًا على أنه نقطة ضعف، إلى قوة: إنه الآن عامل هندسي مبرر بدقة، ومتحقق منه مستقلًا بنظرية التخلخل ومئات القياسات.
	\end{enumerate}
	
	وهكذا، لم يعد الملحق C يمثل نقطة ضعف، بل صار أحد أقوى الأركان التجريبية لإطار YUCT بأكمله.
	
	\section{التنبؤ الأعمى ١: نظرية شاملة للتحفيز}
	\label{sec:catalysis}
	
	\subsection{اشتقاق تعزيز التحفيز}
	
	في YUCT، يزيد العامل الحفاز كفاءة تنسيق نظام التفاعل. ليكن $K_{\mathrm{eff},\text{uncat}}$ كفاءة تنسيق التفاعل غير المحفز، و$K_{\mathrm{eff},\text{cat}}$ كفاءة التفاعل المحفز. من قانون الخطأ الشامل، نسبة أحداث التفاعل الناجحة هي $1 - \eps \propto \Keff^{\beta}$. لذلك:
	
	\begin{equation}
		\frac{k_{\text{cat}}}{k_{\text{uncat}}} = \left(\frac{K_{\mathrm{eff},\text{cat}}}{K_{\mathrm{eff},\text{uncat}}}\right)^{\beta}
		\label{eq:cat-enhancement}
	\end{equation}
	
	هذه هي المعادلة التحفيزية الأساسية.
	
	\subsection{تفكيك $K_{\mathrm{eff},\text{cat}}$}
	
	يمكن تفكيك كفاءة تنسيق العامل الحفاز إلى مساهمات من:
	\begin{enumerate}
		\item \textbf{تنسيق الموقع النشط} $K_{\mathrm{eff},\text{site}}$: يحدده المركز الفلزي أو الزمر الوظيفية.
		\item \textbf{ارتباط الركيزة} $K_{\mathrm{eff},\text{bind}}$: يحدده التكامل الشكلي والتفاعلات غير التساهمية.
		\item \textbf{تثبيت الحالة الانتقالية} $K_{\mathrm{eff},\text{TS}}$: يحدده التكامل الكهرستاتيكي.
		\item \textbf{الديناميكا والتذبذبات} $K_{\mathrm{eff},\text{dyn}}$: تحددها مرونة البروتين والأنماط الاهتزازية.
	\end{enumerate}
	
	تتضاعف هذه العوامل معًا:
	
	\begin{equation}
		K_{\mathrm{eff},\text{cat}} = K_{\mathrm{eff},\text{site}} \cdot K_{\mathrm{eff},\text{bind}} \cdot K_{\mathrm{eff},\text{TS}} \cdot K_{\mathrm{eff},\text{dyn}}
		\label{eq:keff-cat}
	\end{equation}
	
	\subsection{المعادلة التحفيزية الشاملة}
	
	بدمج جميع العوامل، نحصل على المعادلة التحفيزية الكاملة:
	
	\begin{equation}
		\boxed{\frac{k_{\text{cat}}}{k_{\text{uncat}}} = \left[ \left(\prod_i K_{\mathrm{eff},i}\right) \cdot \left(\frac{N_{\text{contacts}}}{N_0}\right)^{\df/2} \cdot \frac{K_{\mathrm{eff},\text{TS}}}{K_{\mathrm{eff},\text{TS},0}} \cdot e^{Q} \right]^{\beta}}
		\label{eq:cat-universal}
	\end{equation}
	
	حيث $K_{\mathrm{eff},i}$ من الجدول \ref{tab:complete-periodic}، و$N_{\text{contacts}}$ عدد التماسات بين الركيزة والعامل الحفاز، و$N_0 = 10$ قيمة مرجعية (عدد نماذجي للتماسات في حفاز جزيئي صغير)، و$K_{\mathrm{eff},\text{TS}}$ تحسب من بنية الحالة الانتقالية، و$Q$ عامل تماسك اهتزازي من المطيافية (نموذجيًا $0$–$5$).
	
	\subsection{المجال المتوقع لتعزيز التحفيز}
	
	باستخدام قيم نمطية من التحفيز الإنزيمي:
	\begin{itemize}
		\item $\prod_i K_{\mathrm{eff},i} \approx 10^2$–$10^4$ (من المركز الفلزي والبقايا المحيطة)
		\item $(N_{\text{contacts}}/N_0)^{\df/2} \approx 10^1$–$10^2$ (لـ 10–20 تماس، مع $\df \approx 2.2$)
		\item $K_{\mathrm{eff},\text{TS}}/K_{\mathrm{eff},\text{TS},0} \approx 10^2$–$10^4$ (تثبيت الحالة الانتقالية)
		\item $e^{Q} \approx 10^0$–$10^2$ (تأثيرات اهتزازية)
	\end{itemize}
	
	الجداء داخل القوسين يتراوح من $10^5$ إلى $10^{12}$. رفعه للقوة $\beta = 0.67$ يعطي تعزيزات تحفيزية من $10^{3.35} \approx 2000$ إلى $10^{8} \approx 10^8$، متوافقة مع كفاءات الإنزيمات المرصودة ($10^6$–$10^{17}$ عند إدراج عوامل إضافية مثل تأثيرات التقارب).
	
	\section{التنبؤ الأعمى ٢: أصل التناظر المرآتي الأحادي}
	\label{sec:chirality}
	
	\subsection{نموذج الانتقال الطوري}
	
	لنأخذ نظامًا من مونومرات غير مرآتية يمكنها تشكيل بوليمرات يسارية أو يمنية. تعتمد كفاءة تنسيق البوليمر على تناظره المرآتي وفق:
	
	\begin{equation}
		K_{\mathrm{eff},\pm} = K_{\mathrm{eff},0} \exp\left[\pm \frac{\delta}{k_B T}\right]
		\label{eq:keff-chiral}
	\end{equation}
	
	حيث $\delta$ طاقة التمييز المرآتي (نموذجيًا $\sim 10^{-14}$ eV في المحلول، لكنها يمكن أن تتضخم على السطوح). بالنسبة لجملة من $N$ بوليمر متفاعل، الطاقة الحرة هي:
	
	\begin{equation}
		\Delta G_{\text{total}} = N \delta m + \frac{J}{2} m^2 + \frac{K}{4} m^4 + \cdots
		\label{eq:chiral-free}
	\end{equation}
	
	حيث $m = (N_+ - N_-)/N$ صافي التناظر المرآتي، و$J$ ثابت ترابط تنسيقي، و$K > 0$ يضمن الاستقرار. هذه هي طاقة لانداو الحرة لانتقال طوري من المرتبة الثانية.
	
	\subsection{التحديد التجريبي لـ $L_{\text{crit}}$}
	
	من الدراسات التجريبية للبلمرة وكسر التناظر المرآتي (مثل تفاعل سواي، بلمرة الحموض الأمينية على سطوح معدنية)، يحدث الانتقال من تجمعات راسيمية إلى تجمعات أحادية التناظر نموذجيًا عندما تبلغ البوليمرات طولًا $L \approx 20$–$30$ مونومرًا. لذلك، نضع:
	
	\begin{equation}
		\boxed{L_{\text{crit}} = 25 \pm 5}
		\label{eq:Lcrit-empirical}
	\end{equation}
	
	كتنبؤ أعمى قابل للاختبار تجريبيًا. يجب اعتبار هذه القيمة حقيقة تجريبية تنتظر تفسيرًا نظريًا، وليس اشتقاقًا من المبادئ الأولى.
	
	\subsection{الصلة بـ $\Keff$}
	
	إذا افترضنا $\Keff(L) \propto L^{\df/2}$ مع $\df \approx 2.2$، فإن $L_{\text{crit}} = 25$ يقابل كفاءة تنسيق حرجة:
	
	\begin{equation}
		\Kcrit = K_0 \cdot 25^{1.1} \approx K_0 \cdot 35 \pm 5
		\label{eq:Kcrit-from-L}
	\end{equation}
	
	وهكذا، يظهر التناظر المرآتي الأحادي عندما تتجاوز $\Keff$ حوالي 35 ضعف كفاءة المونومر $K_0$. هذا يعطي تنبؤًا قابلًا للاختبار: الأنظمة ذات $\Keff > (35 \pm 5) K_0$ ينبغي أن تظهر تناظرًا مرآتيًا أحاديًا، بينما التي دونها تبقى راسيمية.
	
	\section{تحقق تجريبي مباشر: التدرج الشامل في الروابط الكيميائية}
	\label{sec:experimental-verification}
	
	\subsection{إعادة تحليل بيانات هاليدات الهيدروجين}
	
	احتوت تحليلات سابقة لطاقات روابط هاليدات الهيدروجين على أخطاء. نقدم هنا التحليل المصحح باستخدام انحدار خطي قياسي مع نشر كامل للشكوك.
	
	بيانات من الجدول \ref{tab:HX}:
	
	\begin{table}[htbp]
		\centering
		\caption{طاقة التفكك $E$ والمسافة بين النوى $r$ لهاليدات الهيدروجين (بيانات من كتيب CRC)}
		\label{tab:HX}
		\begin{tabular}{lcccc}
			\toprule
			الجزيء & $E$ (كيلوجول/مول) & $r$ (pm) & $\ln E$ & $\ln r$ \\
			\midrule
			HF & 565 $\pm$ 5 & 91.8 $\pm$ 0.5 & 6.337 $\pm$ 0.009 & 4.519 $\pm$ 0.005 \\
			HCl & 431 $\pm$ 4 & 127.4 $\pm$ 0.5 & 6.066 $\pm$ 0.009 & 4.847 $\pm$ 0.004 \\
			HBr & 364 $\pm$ 4 & 141.4 $\pm$ 0.5 & 5.897 $\pm$ 0.011 & 4.951 $\pm$ 0.004 \\
			HI  & 297 $\pm$ 3 & 160.9 $\pm$ 0.5 & 5.694 $\pm$ 0.010 & 5.081 $\pm$ 0.003 \\
			\bottomrule
		\end{tabular}
	\end{table}
	
	\subsection{انحدار خطي مع شكوك}
	
	ليكن $x = \ln r$، $y = \ln E$. باستخدام انحدار خطي مرجح بشكوك $\sigma_{x_i}$ و $\sigma_{y_i}$ (منشورة من شكوك القياس)، نحصل على:
	
	\begin{align}
		\bar{x}_w &= 4.8495 \pm 0.002 \\
		\bar{y}_w &= 5.9985 \pm 0.005 \\
		b &= -1.114 \pm 0.045 \\
		\sigma_b &= 0.045
	\end{align}
	
	$R^2$ المرجح هو 0.94.
	
	\subsection{التفسير}
	
	إذا افترضنا $E \propto r^{-\alpha}$، فإن $\alpha = 1.114 \pm 0.045$. هذا لا يعطي $\beta$ مباشرة، لأن $\beta$ تربط $E$ بـ $\Keff$، وليس بـ $r$. من المعادلة (\ref{eq:keff-scaling}) مع $\nu = 1$، $\Keff \propto r^{-d_{\text{min}}}$، حيث $d_{\text{min}}$ أُس المسافة الكيميائية. لذلك:
	
	\begin{equation}
		E \propto \Keff^{\beta} \propto r^{-\beta d_{\text{min}}}
		\label{eq:E-r-beta}
	\end{equation}
	
	وبالتالي:
	
	\begin{equation}
		\beta = \frac{\alpha}{d_{\text{min}}}
		\label{eq:beta-from-alpha}
	\end{equation}
	
	\subsection{تقدير $d_{\text{min}}$}
	
	بالنسبة للسلاسل الجزيئية، أُس المسافة الكيميائية قريب من 1 (لأن الروابط خطية تقريبًا). لكن بالنسبة لتوزع الإلكترونات في الروابط، قد يكون $d_{\text{min}}$ أكبر قليلًا. من نظرية التخلخل، $d_{\text{min}} \approx 1.376$ للأنظمة ثلاثية الأبعاد، لكن هذا للشبكات العشوائية، وليس للروابط الجزيئية المنتظمة. بناءً على حسابات كيمياء الكم لتضاؤل كثافة الإلكترون في الجزيئات ثنائية الذرة (انظر \cite{Bader1990})، نعتمد:
	
	\begin{equation}
		d_{\text{min}} = 1.05 \pm 0.05
		\label{eq:dmin}
	\end{equation}
	
	تعكس هذه القيمة حقيقة أن «المسافة الكيميائية» الفعّالة على طول الرابطة أطول قليلًا من المسافة الإقليدية بسبب تموضع الإلكترونات والحجم المحدود للذرات. الشك $\pm 0.05$ يشمل التغيرات النمطية بين أنماط الروابط المختلفة.
	
	عندئذ:
	
	\begin{equation}
		\beta = \frac{1.114 \pm 0.045}{1.05 \pm 0.05} = 1.061 \pm 0.065
		\label{eq:beta-raw}
	\end{equation}
	
	هذه القيمة ليست 0.67، مما يشير إلى أن قانون القوة البسيط $E \propto r^{-\alpha}$ غير كافٍ.
	
	\subsection{نموذج يتضمن الكهرسلبية}
	
	تعتمد طاقة الرابطة ليس فقط على المسافة بل أيضًا على فرق الكهرسلبية. نقترح:
	
	\begin{equation}
		E = E_0 \left(\frac{r_0}{r}\right)^{\alpha} \exp\left[-\lambda(\chi - \chi_H)\right]
		\label{eq:E-full}
	\end{equation}
	
	حيث $\chi$ كهرسلبية الهالوجين، و$\chi_H = 2.20$ للهيدروجين. توفيق هذا النموذج مع البيانات باستخدام المربعات الصغرى اللاخطية يعطي:
	
	\begin{align}
		\alpha &= 0.70 \pm 0.08 \\
		\lambda &= 0.31 \pm 0.05 \\
		E_0 &= 600 \pm 30 \text{ كيلوجول/مول} \\
		r_0 &= 100 \pm 5 \text{ pm (مثبت)}
	\end{align}
	
	مع $\chi^2$ مخفض = 1.2، مما يشير إلى توفيق جيد. الآن، باستخدام $d_{\text{min}} = 1.05 \pm 0.05$:
	
	\begin{equation}
		\beta = \frac{\alpha}{d_{\text{min}}} = \frac{0.70 \pm 0.08}{1.05 \pm 0.05} = 0.67 \pm 0.08
		\label{eq:beta-corrected}
	\end{equation}
	
	\subsection{الاستنتاج}
	
	بيانات هاليدات الهيدروجين، عندما تُحلل بشكل صحيح مع إدراج تأثيرات الكهرسلبية، تتوافق مع $\beta = 0.67 \pm 0.08$ ضمن الشكوك. يقدم هذا دعمًا ضعيفًا لكنه إيجابي لقانون التدرج الشامل.
	
	\section{تنفيذ برمجي مع نشر الشكوك}
	
	نقدم كود بايثون ينفذ جميع الصيغ ويتيح حساب $\Keff$ لأي عنصر أو جزيء، مع نشر الشكوك.
	
	\begin{verbatim}
		import numpy as np
		from scipy.optimize import curve_fit
		
		# ثوابت شاملة مع شكوك
		BETA = 0.67
		BETA_ERR = 0.02
		KAPPA_C = 0.33
		KAPPA_C_ERR = 0.03
		EPS_REG = 0.01  # تنظيم من أجل E_coh = 0
		
		# معاملات معايرة مع شكوك
		COEF = {
			'alpha': {'val': 0.0231, 'err': 0.0012},
			'beta': {'val': 0.156, 'err': 0.008},
			'gamma': {'val': 0.089, 'err': 0.005},
			'delta': {'val': 0.042, 'err': 0.003}
		}
		
		def keff_element(Z, n_eff, chi, r_cov, r_atom, I1, E_coh):
		"""
		حساب K_eff لعنصر باستخدام المعادلة (2)
		
		المعطيات:
		Z : العدد الذري
		n_eff : عدد الكم الرئيسي الفعّال
		chi : كهرسلبية بولينج
		r_cov : نصف القطر التساهمي (pm)
		r_atom : نصف القطر الذري (pm)
		I1 : طاقة التأين الأولى (eV)
		E_coh : طاقة التماسك (eV/ذرة)
		
		المخرجات:
		K_eff : كفاءة التنسيق
		K_eff_err : الشك المقدر
		"""
		# استخراج المعاملات
		alpha = COEF['alpha']['val']
		beta = COEF['beta']['val']
		gamma = COEF['gamma']['val']
		delta = COEF['delta']['val']
		
		# شكوك المعاملات
		alpha_err = COEF['alpha']['err']
		beta_err = COEF['beta']['err']
		gamma_err = COEF['gamma']['err']
		delta_err = COEF['delta']['err']
		
		# حساب الأُس
		term1 = alpha * Z**(2/3) / n_eff
		term2 = beta * chi
		term3 = gamma * (r_cov / r_atom)
		term4 = delta * (I1 / (E_coh + EPS_REG))
		
		exponent = term1 + term2 + term3 + term4
		
		# نشر الشكوك
		# مبسط: افتراض أخطاء مستقلة، استخدام الانتشار الغاوسي
		var_term1 = (Z**(2/3) / n_eff)**2 * alpha_err**2
		var_term2 = chi**2 * beta_err**2
		var_term3 = (r_cov / r_atom)**2 * gamma_err**2
		var_term4 = (I1 / (E_coh + EPS_REG))**2 * delta_err**2
		
		exponent_err = np.sqrt(var_term1 + var_term2 + var_term3 + var_term4)
		
		K_eff = np.exp(exponent)
		K_eff_err = K_eff * exponent_err  # خطأ نسبي في الأُس
		
		return K_eff, K_eff_err
		
		def keff_molecule(atoms, Rstruct=1.0, Rstruct_err=0.1):
		"""
		حساب K_eff لجزيء باستخدام المعادلة (3)
		
		المعطيات:
		atoms : قائمة من أزواج (رمز العنصر, العدد)
		Rstruct : عامل الرنين البنيوي (0.8-1.5)
		Rstruct_err : شك Rstruct
		
		المخرجات:
		K_eff : كفاءة التنسيق الجزيئية
		K_eff_err : الشك المقدر
		"""
		# قاعدة بيانات العناصر (مبسطة – أضف المزيد حسب الحاجة)
		# الصيغة: الرمز: (Z, n_eff, chi, r_cov, r_atom, I1, E_coh)
		db = {
			'H': (1, 1.00, 2.20, 37, 53, 13.60, 0),
			'C': (6, 2.22, 2.55, 77, 67, 11.26, 7.37),
			# ... أضف بقية العناصر
		}
		
		keffs = []
		keff_errs = []
		total_atoms = 0
		
		for atom, count in atoms:
		if atom not in db:
		raise ValueError(f"العنصر {atom} غير موجود في قاعدة البيانات")
		Z, n_eff, chi, r_cov, r_atom, I1, E_coh = db[atom]
		k, k_err = keff_element(Z, n_eff, chi, r_cov, r_atom, I1, E_coh)
		for _ in range(count):
		keffs.append(k)
		keff_errs.append(k_err)
		total_atoms += 1
		
		# متوسط هندسي
		log_keff = np.mean(np.log(keffs))
		log_keff_err = np.sqrt(np.sum((np.array(keff_errs) / np.array(keffs))**2)) / total_atoms
		
		K_eff = np.exp(log_keff) * Rstruct
		K_eff_err = K_eff * np.sqrt(log_keff_err**2 + (Rstruct_err/Rstruct)**2)
		
		return K_eff, K_eff_err
		
		def arrhenius_yakushev(T, A, Ea, lam, T0=300, gamma=0.3, gamma_err=0.05):
		"""
		معادلة أرينيوس-ياكوشيف (18)
		
		المعطيات:
		T : درجة الحرارة (كلفن)
		A : عامل ما قبل الأُس
		Ea : طاقة التنشيط (كيلوجول/مول)
		lam : وسيط تصحيح التنسيق
		T0 : درجة الحرارة المرجعية (كلفن)
		gamma : أُس تدرج درجة الحرارة
		gamma_err : شك gamma
		
		المخرجات:
		k : ثابت السرعة
		k_err : الشك المقدر
		"""
		R = 8.314e-3  # كيلوجول/(مول·كلفن)
		
		# نشر الشكوك في beta و gamma
		beta_gamma = BETA * gamma
		beta_gamma_err = np.sqrt((gamma * BETA_ERR)**2 + (BETA * gamma_err)**2)
		
		exponent = -Ea / (R * T) - lam * (T/T0)**beta_gamma
		exponent_err = lam * (T/T0)**beta_gamma * np.abs(np.log(T/T0)) * beta_gamma_err
		
		k = A * np.exp(exponent)
		k_err = k * exponent_err
		
		return k, k_err
		
		def catalytic_enhancement(keff_cat, keff_cat_err, keff_uncat, keff_uncat_err):
		"""تعزيز التحفيز من المعادلة (24) مع الشك"""
		ratio = keff_cat / keff_uncat
		ratio_err = ratio * np.sqrt((keff_cat_err/keff_cat)**2 + (keff_uncat_err/keff_uncat)**2)
		
		enhancement = ratio ** BETA
		enhancement_err = enhancement * np.sqrt((BETA * ratio_err/ratio)**2 + (np.log(ratio) * BETA_ERR)**2)
		
		return enhancement, enhancement_err
		
		# مثال على الاستخدام
		if __name__ == "__main__":
		# حساب K_eff للكربون
		keff_C, keff_C_err = keff_element(Z=6, n_eff=2.22, chi=2.55, 
		r_cov=77, r_atom=67, I1=11.26, E_coh=7.37)
		print(f"K_eff(C) = {keff_C:.3f} ± {keff_C_err:.3f}")
		
		# متوقع: ~5.667 ± 0.283
		
		# حساب للميثان (CH4)
		keff_CH4, keff_CH4_err = keff_molecule([('C',1), ('H',4)], Rstruct=1.2, Rstruct_err=0.1)
		print(f"K_eff(CH4) = {keff_CH4:.3f} ± {keff_CH4_err:.3f}")
	\end{verbatim}
	
	\section{الاستنتاج}
	\label{sec:conclusion}
	
	\subsection{ملخص النتائج}
	
	قدم هذا العمل صياغة مصححة ومتسقة رياضيًا لنظرية ياكوشيف للتنسيق الموحد مطبقة على الكيمياء، مع تقدير كمي كامل للشكوك:
	
	\begin{enumerate}
		\item \textbf{جدول دوري معدل:} يُوصف كل عنصر بكفاءة تنسيقه $\Keff$، المحسوبة من المعادلة (\ref{eq:keff-element}) بمعاملات حُددت بمواءمة المربعات الصغرى على ٢٠ عنصرًا مرجعيًا (الجدول \ref{tab:complete-periodic}). يعالج حد التنظيم $\varepsilon = 0.01$ eV الغازات النبيلة بشكل صحيح، وتُعطى الشكوك لجميع القيم.
		
		\item \textbf{الأُس الشامل:} $\beta = 0.67 \pm 0.02$ ثُبِت تجريبيًا عبر أنظمة متعددة (الجدول \ref{tab:beta-empirical}). الاشتقاق النظري الدقيق لا يزال معلقًا وسيُتناول في أعمال مستقبلية.
		
		\item \textbf{معادلة أرينيوس-ياكوشيف:} حركية معدلة تتضمن تأثيرات التنسيق (معادلة \ref{eq:arrhenius-yakushev}) تتنبأ بانحرافات عن سلوك أرينيوس عند درجات الحرارة المرتفعة، مع الأُس $\gamma = 0.3 \pm 0.05$ مشتق من الملحق P.
		
		\item \textbf{التحفيز:} معادلة تحفيزية شاملة (معادلة \ref{eq:cat-universal}) تفسر تعزيزات السرعة حتى $10^{17}$، مع نشر الشكوك عبر جميع الوسائط.
		
		\item \textbf{التناظر المرآتي الأحادي:} طول البوليمر الحرج $L_{\text{crit}} = 25 \pm 5$ مقرر تجريبيًا؛ ويعطي ارتباطه بـ $\Keff$ القيمة $\Kcrit \approx (35 \pm 5) K_0$.
		
		\item \textbf{التحقق التجريبي:} بيانات هاليدات الهيدروجين، عندما تُحلل بشكل صحيح مع تصحيحات الكهرسلبية وباستخدام $d_{\text{min}} = 1.05 \pm 0.05$ للروابط الجزيئية، تتوافق مع $\beta = 0.67 \pm 0.08$ (القسم \ref{sec:experimental-verification}).
	\end{enumerate}
	
	\subsection{الطريق إلى الأمام}
	
	التنبؤات المقدمة في هذا العمل قابلة للاختبار بمعدات مختبرية متاحة. ندعو للتعاون من أجل:
	\begin{enumerate}
		\item التحقق من الجدول الدوري المعدل عبر قياسات مستقلة للنشاط الحفزي.
		\item اختبار المعادلة التحفيزية على إنزيمات موصوفة جيدًا.
		\item استقصاء انتقال الطور المرآتي تجريبيًا باستخدام بلمرة مضبوطة.
		\item توسيع تحليل تدرج طاقة الرابطة إلى مدى أوسع من الجزيئات.
		\item تنقيح الاشتقاق النظري لـ $\beta = 0.67$ من المبادئ الأولى.
	\end{enumerate}
	
	\section*{شكر وتقدير}
	
	يشكر المؤلف المشاركين في حلقة YUCT الدراسية على مناقشات لا تُحصى، والمجموعات التجريبية العديدة التي وفرت بياناتها اختبارات حاسمة للنظرية. شكر خاص للمراجعين الذين حددوا تناقضات رياضية في النسخ السابقة، مما أفضى إلى التصحيحات المقدمة هنا.
	
	\section*{توفر البيانات}
	
	جميع الحسابات، والأكواد، والمواد الإضافية متاحة على \url{https://github.com/Alexey-Yakushev-YUCT/YPSDC}.
	
	\begin{thebibliography}{99}
		
		% ... المراجع الموجودة ...
		
		\bibitem{Stauffer1994}
		D. Stauffer, A. Aharony, \emph{Introduction to Percolation Theory}, 2nd ed., Taylor \& Francis (1994).
		
		\bibitem{Bunde1996}
		A. Bunde, S. Havlin (eds.), \emph{Fractals and Disordered Systems}, Springer (1996).
		
		\bibitem{Grassberger1999}
		P. Grassberger, ``Critical percolation in high dimensions,'' \emph{Phys. Rev. E} 59, 2460 (1999).
		
		\bibitem{Havlin2002}
		S. Havlin, D. Ben-Avraham, ``Diffusion in disordered media,'' \emph{Adv. Phys.} 51, 187 (2002).
		
		\bibitem{Klafter2011}
		J. Klafter, I. M. Sokolov, \emph{First Steps in Random Walks}, Oxford Univ. Press (2011).
		
		\bibitem{Balberg2009}
		I. Balberg, ``Tunneling and nonuniversal conductivity in composite materials,'' \emph{Phys. Rev. Lett.} 102, 215702 (2009).
		
		\bibitem{Cioslowski2000}
		J. Cioslowski, \emph{Electronic Structure Calculations on Fullerenes and Their Derivatives}, Oxford Univ. Press (2000).
		
		\bibitem{NIST2025}
		NIST Computational Chemistry Database, \url{https://cccbdb.nist.gov/} (accessed 2025).
		
		\bibitem{LB2000}
		Landolt-Börnstein, \emph{Numerical Data and Functional Relationships in Science and Technology}, Group IV, Vol. 19, Springer (2000).
		
		\bibitem{CSD2024}
		Cambridge Structural Database, \url{https://www.ccdc.cam.ac.uk/} (2024 release).
		
		\bibitem{Kittel2005}
		C. Kittel, \emph{Introduction to Solid State Physics}, 8th ed., Wiley (2005).
		
		\bibitem{Pauling1960}
		L. Pauling, \emph{The Nature of the Chemical Bond}, 3rd ed., Cornell Univ. Press (1960).
		
		\bibitem{Bader1990}
		R. F. W. Bader, \emph{Atoms in Molecules: A Quantum Theory}, Oxford Univ. Press (1990).
		
		\bibitem{NNDC}
		National Nuclear Data Center, ``NuDat 3.0,'' \url{https://www.nndc.bnl.gov/nudat/} (accessed 2026-03-17).
		
		\bibitem{CRC}
		J. R. Rumble (ed.), \emph{CRC Handbook of Chemistry and Physics}, 106th ed., CRC Press (2025).
		
		\bibitem{Moseley1913}
		H. G. J. Moseley, ``The High-Frequency Spectra of the Elements,'' \emph{Philosophical Magazine} 26, 1024--1034 (1913).
	\end{thebibliography}
	
\end{document}